\section{C-Tree Math}\label{sec:v8-ctree-math}

The next three sections give the math. This one defines the C-Tree
objects and the three local laws.
Section~\ref{sec:v8-objective} turns the laws into a training
objective. Section~\ref{sec:v8-audit-certificate} makes the audit
operational. Throughout, we work over a fixed manifesto and a fixed
realized tree.

\subsection{Objects and States}

Let \(\mathcal X\) be the input domain for manifestos and manifesto
spans, let \(\mathcal Y\) be the task-output space, and let
\(\mathcal Z\) be the policy-state space. In the manifesto LLM
application, \(\mathcal X\) and \(\mathcal Z\) are both spaces of text:
\(\mathcal X\) is raw manifesto text and \(\mathcal Z\) is short
policy-evidence summaries written by \(g\). \(\mathcal Y\) is the
seven-point policy score on one dimension, or the vector of six
dimension scores when one tree serves all six.

Fix a manifesto \(x\in\mathcal X\). A realized C-Tree is a finite
ordered binary tree \(T\) with node set \(V(T)\) and leaves forming a
fixed contiguous partition of \(x\). Each node \(v\in V(T)\) represents
a span \(x_v\in\mathcal X\): at a leaf, \(x_v\) is the raw manifesto
span; at an internal node with left child \(u\) and right child \(w\),
\(x_v=x_u\concat x_w\). Here \(\concat\) denotes ordered concatenation.
The root span is \(x_{\mathrm{root}}=x\).

The C-Tree learns one state map, written \(g\). At a leaf, \(g\) reads
a raw manifesto span and returns a policy-evidence state. At an
internal node, \(g\) reads the two child states and returns the parent
state. The same \(g\) can also be applied to a state that \(g\) already
produced, e.g.\ during a cleanup pass that rewrites a state before it
merges with its sibling. The stored state at node \(v\) is
\[
  z_v =
  \begin{cases}
    g(x_v), & v \text{ is a leaf},\\
    g(z_u,z_w), & v \text{ has children }u,w.
  \end{cases}
\]

A scorer reads one state and returns a number. We write \(f(z)\) for a
scorer applied to state \(z\), \(f^\ast\) for the target oracle (the
expert or rubric scorer for one policy dimension), and
\(\widehat f\) for the fitted cheaper proxy used for dense local
supervision. In the manifesto LLM application, the scorer is a prompt
that places the state where it expects the document, and the proxy
\(\widehat f\) is the same prompt applied to a smaller LLM. \(f^\ast\)
is observed only when the oracle is called.

For any scorer \(q\), query equivalence is
\[
  a\sim_q b
  \quad\Longleftrightarrow\quad
  d_Y(q(a),q(b))=0 .
\]
Here \(d_Y\) is the task-space discrepancy on \(\mathcal Y\), and
\(a\sim_q b\) means that \(a\) and \(b\) give the same answer under
\(q\). We specialize to the oracle query \(q=f^\ast\) and write
\[
  a\sim_\ast b
  \quad\Longleftrightarrow\quad
  d_Y(f^\ast(a),f^\ast(b))=0 .
\]
This convention treats an \(f^\ast\) call as observation-error-free
once it is made; the estimation problem is that such calls are sparse.
Section~\ref{sec:v8-audit-certificate} replaces the equality with a
small tolerance and bounds the resulting drift in the reported number.

\begin{table}[t]
  \centering
  \small
  \begin{tabular}{@{}p{0.22\linewidth}p{0.68\linewidth}@{}}
    \toprule
    Symbol & Meaning \\
    \midrule
    \(x,T,V(T),v\) & Manifesto, realized C-Tree, node set of the tree, and a node. \\
    \(u,w\) & Left and right child nodes of an internal node. \\
    \(x_v,z_v\) & Manifesto span represented by node \(v\), and the policy state stored at \(v\). \\
    \(g\) & One learned C-Tree state map/program, called on raw leaves, on already-produced states, and on pairs of child states. \\
    \(\mathcal X,\mathcal Y,\mathcal Z,\concat\) & Input domain, task-output space, state space, and ordered concatenation. \\
    \(d_Y,\sim_\ast\) & Task-space discrepancy and oracle equivalence. \\
    \(f^\ast,\widehat f\) & Target oracle readout and fitted proxy readout. \\
    C1/C2/C3 & Sufficiency, idempotence, and merge consistency (the local laws). \\
    \bottomrule
  \end{tabular}
  \caption{\textbf{Core notation for the C-Tree math.}}
  \label{tab:v8-core-notation}
\end{table}

\begin{assumption}[Fixed Realized Tree]\label{ass:v8-fixed-tree}
For each document \(x\), the leaf partition and tree topology are conditioned
on as fixed when stating the local laws and the objective.
\end{assumption}

The assumption isolates representation validity from partition
selection. A deterministic partition rule returns the same statement
after conditioning on its realized tree. The leaf partition is a
chunker design choice; in the manifesto experiments it is a fixed
token budget per leaf, swept from \(256\) to \(8{,}192\) tokens
(Section~\ref{sec:v8-benoit-tree-setup}).

\subsection{Local Laws}

The three local laws are stated on the realized calls made by the tree.

\begin{definition}[Oracle C-Tree Local Laws]\label{def:v8-local-laws}
For a realized tree \(T\), state map \(g\), and target oracle \(f^\ast\), the
local laws are:
\[
\begin{aligned}
\mathrm{C1:}\quad&
g(x_v)\sim_\ast x_v
&&\text{for every leaf }v,\\
\mathrm{C2:}\quad&
g(z_v)\sim_\ast z_v
&&\text{for every produced state }z_v,\\
\mathrm{C3:}\quad&
g(z_u,z_w)\sim_\ast x_u\concat x_w
&&\text{for every internal node }v=(u,w).
\end{aligned}
\]
\end{definition}

C1 (sufficiency) says that once a raw manifesto span enters the tree,
the state \(g(x_v)\) is already sufficient for oracle-facing questions
about that span. A leaf state may be shorter, more structured, or
more readable than the raw text, but it must stay in the same
\(f^\ast\)-fiber.

C2 (idempotence) says that summarizing a state again gives the same
expert-facing answer as the state itself: a summary of a summary
lands in the same \(f^\ast\)-fiber as the summary. This prevents a
cleanup pass from turning a conditional policy claim into a different
expert-facing position.

C3 (merge consistency) says that merging two child states produces a
parent state equivalent to the represented parent span
\(x_u\concat x_w\). This is where cross-section policy meaning
matters. The parent state must preserve interactions that are
invisible from scalar child scores: qualifications, contradictions,
cross-references, and tradeoffs between tax, spending, welfare,
environmental, or decentralization claims.

The local laws define which learned state maps count as sensible for
the task; they double as post hoc diagnostics. The exact statement
above assumes \(f^\ast\) can be called at every node, which is too
expensive in deployment. So dense local training uses the proxy
\(\widehat f\) at every node, and a small set of sampled \(f^\ast\)
calls then estimates and corrects the gap between the proxy and the
oracle. Section~\ref{sec:v8-objective} writes the objective for this
setup; Section~\ref{sec:v8-audit-certificate} writes the correction.

\subsection{Why States Matter}

The local laws are state laws, not scalar-score laws. If an economic-policy
leaf stores only \(f^\ast(x_v)\), then the parent has no access to evidence
that may matter only after two spans are read together. Two child spans can
have the same local score and still differ in whether they supply a tax
qualification, spending condition, or regional-autonomy exception. The state
\(z_v\) is the carrier for that interaction information.

This is the sense in which a C-Tree can support more than one downstream
readout of the same compressed manifesto, provided those readouts
respect the oracle equivalence relation. The theorem below states that
if the states satisfy the local laws, then the root state lies in the
same \(\sim_\ast\) class as the original manifesto. Anything that is
constant on those classes can then be evaluated at the root.

\begin{theorem}[Local Laws Preserve the Root]\label{thm:v8-root-preservation}
If C1, C2, and C3 hold exactly for \(g\), \(T\), and \(f^\ast\), then the root
state preserves the document oracle:
\[
  d_Y\!\left(f^\ast(z_{\mathrm{root}}),f^\ast(x)\right)=0 .
\]
For stochastic C-Tree reductions and \(m\ge 1\) rounds of
re-summarizing produced states, the expected oracle distortion is
zero.
\end{theorem}

\begin{proof}[Proof sketch]
Induction on the tree. C1 places each leaf state in the
\(f^\ast\)-fiber of its raw span. C3 propagates fiber equivalence
across each merge: \(g(z_u,z_w)\sim_\ast x_u\concat x_w\) once the
child states are fiber-equivalent to their raw spans. C2 absorbs the
case where a state is summarized again before it merges. At the root,
\(z_{\mathrm{root}}\sim_\ast x\). The stochastic version takes
expectations over the realized reductions. Full proof in
Appendix~\ref{app:v8-full-proofs}; machine-checked as
\texttt{one\_pass} (and the multi-round generalization
\texttt{multi\_round\_proper}) in
\texttt{FormalProofs/OPT/PreservationTheorems.lean} and
\texttt{FormalProofs/OPT/ExpectationTheory.lean}.
\end{proof}

\begin{corollary}[Oracle-Compatible Readouts Pass Through C-Tree Nodes]
\label{cor:v8-preferences-pass-to-root}
Call a readout \(\Phi\) \emph{oracle-compatible} if
\[
  a\sim_\ast b \quad\Longrightarrow\quad \Phi(a)=\Phi(b).
\]
Under exact C1/C2/C3 local laws, every realized node \(v\in V(T)\) satisfies
\[
  \Phi(z_v)=\Phi(x_v).
\]
Taking \(v=\mathrm{root}\) gives the root statement
\(\Phi(z_{\mathrm{root}})=\Phi(x)\).

For pairwise preferences, let \(P\) be a comparison rule that is
oracle-compatible in both arguments:
\[
  a\sim_\ast a',\ b\sim_\ast b'
  \quad\Longrightarrow\quad
  P(a,b)=P(a',b').
\]
Then any two valid node states, from the same C-Tree or from different
C-Trees satisfying the same oracle laws, can replace their raw blocks:
\[
  P(z_v,z_{v'})=P(x_v,x_{v'}).
\]
\end{corollary}

\begin{proof}[Proof sketch]
At each node \(v\), apply Theorem~\ref{thm:v8-root-preservation} to
the subtree rooted at \(v\) to get \(z_v\sim_\ast x_v\).
Oracle-compatibility of \(\Phi\) then gives
\(\Phi(z_v)=\Phi(x_v)\). The pairwise variant uses
oracle-compatibility separately in each argument. Full proof in
Appendix~\ref{app:v8-full-proofs}; machine-checked as
\texttt{paper\_preference\_stack\_same\_argmin} in
\texttt{FormalProofs/OPT/MainTheorems.lean}, with method-specific
surfaces \texttt{dpo\_equivalence}, \texttt{grpo\_equivalence}, and
\texttt{grpo\_rl\_equivalence} in the same module.
\end{proof}

The node-level form is what makes the local-law losses useful as local
supervision rather than only as an indirect route to a final document score. If
the raw block \(x_v\) is a meaningful unit for a task-facing readout, then a
valid state \(z_v\) is meaningful for the same readout, provided the readout
depends only on the oracle equivalence class. The root is only the largest
such block. The same applies across documents: two states from
different manifestos can replace their raw spans for any comparison
that depends on the policy oracle, as long as both trees satisfy the
oracle laws on the same dimension.

The restriction is essential. A preference about information not represented in
\(f^\ast\) is not certified by this theorem, even if the root preserves
\(f^\ast\) exactly. To certify that preference, the oracle must be expanded to
include the relevant readout, or the local laws must be stated for a richer
query family. Within the chosen oracle interface, the corollary is the
preference-learning reason the exact theorem matters: a valid C-Tree root can
stand in for the original long document for scores that respect that interface.

This is the sense in which the policy-state story connects to
classical mergeable summaries (Section~\ref{sec:v8-scope-related}
develops the connection). A classical sketch chooses the state by
hand and proves the merge law. A learned C-Tree chooses a state class
and uses local-law losses plus audits to decide whether the fitted
state behaves as a valid task-relative state machine.

\subsection{From Exact Laws to Learning}

The exact statement assumes the oracle is available wherever the
analyst needs it. If \(f^\ast\) could be called at every leaf, state,
and merge, the learning problem would be almost direct: search for a
\(g\) whose C1/C2/C3 discrepancies are small under \(f^\ast\), while
also fitting the document-level expert-survey target.

In the manifesto LLM setting, oracle calls are not free at every node.
An \(f^\ast\) call may be an expert label, an expensive teacher
judgment, or a local policy-evidence label derived from
application-aligned coding. The oracle local-law loss is therefore
observable at only a small set of sampled nodes. C-TreePO optimizes
one learned state map \(g\) for two purposes at once: it should be
useful at the root, and it should behave like a valid C-Tree state
throughout the realized tree.

A realized tree makes this approximation problem concrete. At a leaf,
we ask whether the state preserves the raw policy evidence. At a
produced state, we ask whether summarizing it again changes the
expert-facing content. At an internal node, we ask whether merging
child states preserves the concatenated span. These are the same
semantic requirements used in the exact theorem.
Section~\ref{sec:v8-objective} turns them into an optimization target,
and Section~\ref{sec:v8-audit-certificate} explains how sparse oracle
calls estimate the local-law channel when dense training uses
\(\widehat f\).
